\subsection{多域信号处理}
\label{sec:multi-domain}

单域信号表示存在固有的权衡：时域分析保留时间顺序但忽略频谱内容，而频域分析揭示频谱分量却以时间顺序为代价 \cite{yang2023unsupervised}。这一限制对于激光超声信号尤为重要，因为到达时间和频率内容都携带诊断相关信息。

多域融合通过从时域、频域和时频表示中联合提取特征来解决这一问题 \cite{perry2026health}。时频变换如短时傅里叶变换(STFT)、连续小波变换(CWT)和小波包分解提供互补视角：STFT提供均匀的时频分辨率，而小波使其分辨率适应频率 \cite{chen2019wavelet}。Perry等人证明，融合时域、频域和时频域可改善复合材料中导波的 健康指标提取 \cite{perry2026health}。

尽管有此已验证的益处，多域融合在激光超声去噪中仍未得到充分探索。现有的深度学习方法主要在单一域（时域或频域）中运行，无法联合捕获跨表示的互补信息。这一差距促使本文提出多域融合策略。
